\chapter{Resultados: Dashboards e Indicadores}
\label{cap:resultados}

\comment[inline]{Apresentar o produto final: os dashboards construídos e os indicadores calculados. Mostrar capturas de tela e descrever o que cada visualização revela. Aproximadamente 8--12 páginas.}

\section{Visão Geral da Execução Orçamentária}

\comment[inline]{Dashboard de visão geral: total empenhado vs. liquidado vs. pago por ano, evolução temporal, percentual de execução. Captura de tela + análise do que o painel permite observar que não seria possível no portal original.}

\section{Execução por Função de Governo}

\comment[inline]{Distribuição de despesas por função (educação, saúde, infraestrutura, etc.). Comparação entre anos. Permitir responder: ``Quanto foi investido em educação nos últimos três anos?''. Captura + análise.}

\section{Gasto Per Capita e Comparações}

\comment[inline]{Indicadores normalizados por habitante (usando dados populacionais do IBGE). Comparações com outros municípios de porte similar se disponível. Discutir relevância para o controle social.}

\section{Concentração de Fornecedores}

\comment[inline]{Análise dos principais favorecidos por área/órgão. Índice de concentração (ex.: participação dos top-10 fornecedores no total pago). Permite detectar padrões anômalos de concentração. Captura + análise.}

\section{Receitas Municipais}

\comment[inline]{Se os dados de receita foram incluídos no DW: composição da receita (IPTU, ISS, transferências, etc.), evolução temporal, comparação com despesas (equilíbrio fiscal). Captura + análise.}

\section{Análise dos Dados de Juiz de Fora}

\comment[inline]{Seção analítica: o que os dados revelam sobre a execução orçamentária de Juiz de Fora no período analisado? Não apenas descrever os painéis, mas interpretar os resultados. Quais padrões chamam atenção? O que um cidadão poderia fazer com essa informação? Esta seção é onde a perspectiva freireana se materializa.}
